\documentclass[12pt,twoside]{article}
\usepackage{fancyhdr,makeidx,times,multicol,geometry}
\usepackage{keyval,ifthen,mparhack,manyfoot,poemscol}
\newcounter{scounter}
\newcommand\step{\stepcounter{scounter}\Roman{scounter}}

\title{Postal del Cielo}
\author{J.J. Gómez Cadenas}
\begin{document}
\maketitle

\poemtitle{Adiós a París}
\begin{poem}
\begin{stanza}
¿Debo entonces hacer como Antonio?\\
¿Abrir las ventanas de la senectud que llega\\
y escuchar cómo se aleja la tropa divina,\\
con exquisita música y risas y danzas?
\end{stanza}
\begin{stanza}
¿Debo despedirme de los hoteles de París,\\
del hermoso antro de nuestra primera noche,\\
y del burdel, todo alfombras y espejos, alto en Montmartre,\\
del hotel arruinado, glorioso en su decadencia,\\
donde un viejo camarero nos preguntaba cada mañana,\\
si habíamos disfrutado de nuestro croissant?
\end{stanza}
\begin{stanza}
¿Debo despedirme\\
de las habitaciones, llenas de trastos y mochilas,\\
donde dormíamos los cuatro,\\
---y de las risas, vuestras peleas de niños,\\
la forma en que os colabais en nuestra cama de madrugada---,\\
en aquellos años donde el tiempo no pasaba?
\end{stanza}
\begin{stanza}
¿Debo decirle adiós a vuestros rostros de entonces,\\
a los eternos porqués, a las historias de ruta?\\
¿Decirle adiós a la música que sonaba\\
en nuestro coche atiborrado, maletas y bártulos,\\
papá y mamá, el periquito y los niños,\\
ah, y el dios Baco acompañándonos?
\end{stanza}
\begin{stanza}
El dios Baco que caminaba a nuestro lado,\\
por las calles de una ciudad siempre nueva.
\end{stanza}
\begin{stanza}
¿Debo decir adiós? ¿Eso es todo?\\
¿Despedirme de París,\\
como Antonio se despidió de Alejandría?\\
¿Aceptar, con valor y sin amargura, que el dios se marcha?
\end{stanza}
\begin{stanza}
O debo decir que no\\
y asirme, como un náufrago se aferra a un madero que flota,\\
a la luz que aún queda del día.
\end{stanza}
\end{poem}

\end{document}
